\section{Related Work}
\label{sec:related}

This section reviews the literature on the components of the pipeline
in Figure~\ref{fig:architecture} and on the modelling problem this
article addresses: volumetric video delivery and the transports
proposed for it, subjective assessment of point-cloud video and of
temporal smoothness, and parametric \gls{qoe} models. The literature
was assembled in three passes: a seed set from the bibliographies of
the two conference studies this article extends and of the tutorial by
van der Hooft et al.~\cite{Immersivevideo}, followed by backward
snowballing; the identification of the closest prior model, Udora et
al.~\cite{qoeholographic,udora2026tmm}, and an examination of the works
citing it; and keyword searches of IEEE Xplore, the ACM Digital Library
and Google Scholar for 2023--2026 publications on the same topics.
Only peer-reviewed publications and standards were retained.
Table~\ref{tab:positioning} positions \crystal against the closest
studies.

\begin{table}[!t]
  \caption{Positioning of \crystal against related work.
  \yes{} = supported, \no{} = not addressed, \partially{} = partial.
  ``Network impairment'': delay and loss as controlled independent
  variables (\partially{} denotes bandwidth-only adaptation or offline
  loss injection); ``Reliable \quic'': unmodified, fully reliable \quic
  (\partially{} denotes reliable \gls{tcp} or a modified \quic);
  ``Mechanistic model'': the explanatory variable is derived from
  transport mechanics rather than fitted to delay and loss directly.}
  \label{tab:positioning}
  \centering
  \footnotesize
  \setlength{\tabcolsep}{3pt}
  \renewcommand{\arraystretch}{1.15}
  \begin{tabular*}{\textwidth}{@{\extracolsep{\fill}}llcccccc@{}}
    \toprule
    \thead[l]{Work} & \thead[l]{Codec} &
    \thead{Subjective\\Study} & \thead{VR\\Headset} &
    \thead{Network\\Impairment} & \thead{QoE\\Model} &
    \thead{Reliable\\\quic} & \thead{Mechanistic\\Model} \\
    \midrule
    van der Hooft et al.~\cite{vanderhooft2020qomex} & \vpcc  & \yes & \no        & \partially & \no  & \no        & \no \\
    Subramanyam et al.~\cite{subramanyam2022} & CWI-PCL    & \yes & \yes       & \partially & \no  & \no        & \no \\
    Cox et al.~\cite{Cox2023}               & \vpcc        & \yes & \no        & \no        & \no  & \no        & \no \\
    Nguyen et al.~\cite{nguyen2023}         & \vpcc        & \yes & \partially & \no        & \no  & \no        & \no \\
    Zhou et al.~\cite{zhou2025tomm}         & \vpcc, \glsentryshort{gpcc} & \yes & \yes       & \no        & \no  & \no        & \no \\
    Van Damme et al.~\cite{vandamme2024}    & \vpcc        & \yes & \no        & \no        & \yes & \no        & \no \\
    Dumic \& da Silva Cruz~\cite{dumic2023} & \vpcc        & \yes & \no        & \partially & \no  & \no        & \no \\
    Ravuri et al.~\cite{ravuri2023}         & \vpcc        & \no  & \no        & \yes       & \no  & \partially & \no \\
    Cetus~\cite{hou2026cetus}               & --           & \no  & \no        & \yes       & \no  & \partially & \no \\
    Sidhu \& Bentaleb~\cite{gpcc2026quic}   & \glsentryshort{gpcc}        & \no  & \no        & \partially & \no  & \yes       & \no \\
    Rudolph et al.~\cite{rudolph2025rabbit} & \vpcc        & \no  & \no        & \no        & \no  & \no        & \no \\
    Udora et al.~\cite{qoeholographic,udora2026tmm} & uncompressed & \yes & \yes & \yes & \yes & \partially & \no \\
    \mosaic~\cite{mosaic2026}               & \vpcc        & \yes & \yes       & \yes       & \no  & \yes       & \no \\
    \vquest~\cite{vquest2026}               & frame-indep. & \yes & \yes       & \yes       & \no  & \yes       & \no \\
    \midrule
    \textbf{\crystal (ours)}                & \vpcc        & \yes & \yes       & \yes       & \yes & \yes       & \yes \\
    \bottomrule
  \end{tabular*}
\end{table}

\subsection{Volumetric Video Delivery and Transport}
\label{sec:rw-systems}

Surveys of immersive media delivery describe the pipeline from capture
to \gls{6dof} rendering and identify \gls{qoe} under network impairment
as an open problem~\cite{Immersivevideo,Viola2023,shen2023toward};
Liu et al.~\cite{pcstreaming21} note that measuring the \gls{qoe} of point-cloud streaming is still in
its infancy, and Mozhganfar et al.~\cite{mozhganfar2026noms} classify
the \gls{qoe} models that do exist for volumetric streaming by input
factors and modelling strategy. The standards bodies have catalogued
the influence factors without modelling them: ITU-T
G.1035~\cite{itu_g1035} lists delay, bandwidth and loss among the
network factors and frame rate among the media factors that shape
\gls{vr} \gls{qoe}, and ITU-T P.1320~\cite{itu_p1320} extends the
catalogue to \gls{xr} meetings; neither maps those factors to a score
for volumetric content.

Streaming systems for \vpcc address bandwidth and adaptation. Rudolph
et al.~\cite{rudolph2025rabbit} transcode \vpcc streams in real time by
re-encoding only the video sub-streams; Shi et al.~\cite{qv424} add
viewpoint-aware tiling to \vpcc using the encoder's intermediate
outputs, with a \gls{qoe} model trained on VOLVQAD ratings.
Codec-agnostic systems exploit inter-frame correlation~\cite{hermes}
or neural super-resolution~\cite{yuzu}. These works evaluate with objective metrics or
simulation and hold the encoder configuration, in particular the \gof
length, fixed.

On the transport side, the dominant response to loss is to relax
reliability. Ravuri et al.~\cite{ravuri2023} integrate \gls{dash} with
partially reliable \quic for \vpcc, deciding per object whether to
deliver reliably, and report 46\% higher throughput and 33\% fewer
playout interruptions than reliable \quic at 5\% loss on 5G traces,
with a headless player. Haseeb et al.~\cite{haseeb2025} combine
adaptive encoding with a mixed-reliability \quic protocol; Freeman et
al.~\cite{moq2025} obtain implicit server-side adaptation from
Media-over-\quic delivery timeouts; and Hou et al.~\cite{hou2026cetus}
couple an application-level frame-drop policy to partially reliable
transmission for live volumetric video, reporting frame-rate and
objective-quality gains. A second response handles loss below the
transport: Qian et al.\ duplicate packets in eBPF between network nodes
so that volumetric frames meet their deadlines despite path
instability~\cite{qian2024ebpf}, and Thomas and
Xylomenos~\cite{thomas2025} stream point clouds with forward error
correction on a real 5G testbed without \quic. The one recent pipeline
that keeps production \quic unmodified, Sidhu and Bentaleb's
\gls{dash}-over-\quic delivery of \gpcc~\cite{gpcc2026quic}, finds that
\quic implementation choices materially affect delivery but evaluates
objectively, without a headset or a subjective study. Each of these
works modifies the fully reliable protocol that production stacks
ship~\cite{rfc9000}, handles loss outside it, or omits perceptual
evaluation; that reliable \quic turns loss into delay, and delay into
\gof stalls, is noted in several of them, but its perceptual cost has
not been quantified. Before~\cite{mosaic2026}, no study combined real
\vpcc transmission over \quic, playback in a \gls{vr} headset,
controlled impairment on a real link, 5G delivery and a standardised
subjective study, and no study yet combines those with a model.

\subsection{Subjective Quality Assessment}
\label{sec:rw-subjective}

Subjective studies of point-cloud video have concentrated on
compression and bandwidth. Van der Hooft et al.~\cite{vanderhooft2020qomex}
conduct the first subjective evaluation of adaptive \vpcc streaming,
with 30 participants rating pre-rendered viewports on a flat screen
under bandwidth and buffer variation at a fixed latency. Subramanyam
et al.~\cite{subramanyam2022} evaluate tiled user-adaptive streaming
for \gls{vr} remote communication with 33 users in head-mounted
displays, varying bandwidth. Cox et al.~\cite{Cox2023} release VOLVQAD,
7,680 ratings of 376 \vpcc sequences from 120 participants on 2D
displays, and find texture compression more damaging than geometry.
Nguyen et al.~\cite{nguyen2023} characterise \gls{qoe} and immersion of
\vpcc point clouds on the HoloLens~2 with 32 participants across
quantisation, quality switches, viewing distance and content. Zhou et
al.~\cite{zhou2025tomm} collect quality scores and eye-tracking maps
from 40 participants viewing \vpcc, \gpcc and CWI-PCL content with
\gls{6dof} in \gls{vr}. Alexiou et al.~\cite{Alexiou2023} review the
methodology; Ak et al.~\cite{ak2024basics} benchmark objective metrics
on a broad static corpus, where \pcqm~\cite{pcqm} is among the most
stable across codecs.

Few studies vary the network. Dumic and da Silva Cruz~\cite{dumic2023}
inject \SIrange{0.5}{2}{\percent} packet loss into \vpcc sub-bitstreams
before decoding and report \mos drops exceeding one scale unit at
\SI{0.5}{\percent}, with geometry and occupancy losses far more
damaging than attribute losses. Loss is applied to the bitstream
rather than produced by a transport, so the stimulus is
corruption rather than delay, and the content is shown as 2D video.
Our earlier study~\cite{mosaic2026} rated \vpcc in head-mounted
displays under real \quic impairment but paired each network tier with
a different encoder configuration, so the two factors could not be
separated. The only other subjective study under delay and loss on
reliable transport in head-mounted displays is that of Udora et al.,
discussed in Section~\ref{sec:rw-models}.

The attribute rated in this article is temporal smoothness, which has
a perceptual literature of its own outside point clouds. Judder on
head-mounted displays degrades the perceived quality of
omnidirectional video and interacts with compression
level~\cite{mahmoudpour2019judder}, and Klein et
al.~\cite{klein2024vft} show that frame-timing variations of about
\SI{12}{\milli\second} reduce perceived smoothness while variations of
about \SI{4}{\milli\second} barely do. These studies establish
smoothness as a ratable attribute distinct from fidelity; they do not
connect it to transport behaviour.

\subsection{\gls{qoe} Models}
\label{sec:rw-models}

Parametric \gls{qoe} models relate an impairment to a quality score.
The exponential form used in \crystal descends from the IQX hypothesis
of Fiedler et al.~\cite{fiedler2010}, which derives
$\mathrm{QoE} = a\,e^{-bx} + c$ for a single impairment $x$; Reichl et
al.~\cite{reichl2010} argue for a logarithmic law, and the two agree in
the operating region of interest~\cite{alreshoodi2013}. For buffered
video the impairment that matters is stalling: Ho{\ss}feld et
al.~\cite{hossfeld2011youtube} find an exponential dependence of \mos
on the number and length of stalling events, and the survey by Seufert et
al.~\cite{seufert2015survey} consolidates the stalling ratio as the
established temporal \gls{qoe} variable for HTTP streaming. ITU-T
P.1203~\cite{itu_p1203} standardises that view, folding initial delay,
rebuffering and quality switches into a session score. In all of these
the stalling term is \emph{measured} at the client; \crystal keeps the
exponential kernel and a stalling-type variable but \emph{derives} it
from transport mechanics, so that delay, loss and \gof length enter
only through it.

For volumetric media, three kinds of model have been proposed.
Streaming optimisers such as QV4~\cite{qv424} and Wang et
al.~\cite{wang2022tmm} embed a weighted sum of quality, quality
fluctuation and rebuffering as their objective; these are calibrated
on 2D-video conventions and evaluated in simulation, and Nguyen et
al.~\cite{qoedynamicpc} adapt P.1203 to dynamic point clouds in
\gls{ar} with the parameters of the 2D standard. Feature-based models
map content or device features to quality without a network variable:
Van Damme et al.~\cite{vandamme2024} train a
clustering-based psychometric model on no-reference image features for
\vpcc video at video, \gof and frame granularity, and Zhang et
al.~\cite{zhang2019towards} build a fuzzy inference system for the
HoloLens.

The third kind, and the closest to this work, is the network-aware
regression of Udora et al.~\cite{qoeholographic,udora2026tmm} for live
volumetric telepresence with uncompressed per-frame point clouds over
reliable \glsentryshort{tcp}. From a double-stimulus study with 25
participants in head-mounted displays they fit
\begin{equation}
  \mathrm{QoE} = a\,e^{b D} + c\,\plr + d,
  \label{eq:udora}
\end{equation}
with the coefficients re-estimated for each of three point-cloud
density levels and held-out \gls{plcc} of 0.57--0.66; the journal
version~\cite{udora2026tmm} embeds the model in a learning-based scheme
that adapts resolution and frame rate to user intent and network
state, and leaves \quic to future work. Two features make it the right
baseline. It is the only published model that predicts a subjective
score for volumetric content from network conditions under reliable
transport, so its data carry the same loss-to-delay signature that
\crystal sets out to explain; and its per-frame delivery is the
limiting case of no \gof coupling, which gives the \gof term of
\crystal an external reference point. \crystal differs in form (one
perceptual parameter instead of per-level coefficients), in
explanation (the delay--loss interaction is derived rather than
fitted, with its additive form treated as a prediction to be tested)
and in transport (production \quic rather than \glsentryshort{tcp}).
